\section{Transformer-Based Semantic Segmentation}

\subsection{Principle and Methodology}

Transformer architecture was originally proposed for sequence modeling in natural language processing and later adapted to computer vision tasks \cite{vaswani2017attention}. 
The key component of a transformer is the self-attention mechanism, which enables the model to capture long-range dependencies between input elements.

Given an input feature sequence $X \in \mathbb{R}^{N \times d}$, the attention mechanism computes the relationship between tokens using query ($Q$), key ($K$), and value ($V$) matrices:

\begin{equation}
Q = XW_Q,\quad K = XW_K,\quad V = XW_V
\end{equation}

The scaled dot-product attention is defined as:

\begin{equation}
Attention(Q,K,V) = softmax\left(\frac{QK^T}{\sqrt{d_k}}\right)V
\end{equation}

where $d_k$ represents the dimension of the key vectors used for normalization. 
This mechanism allows each token to attend to all other tokens in the sequence, enabling the model to capture global contextual information \cite{vaswani2017attention}.

In computer vision, images are first divided into patches before being processed by the transformer architecture. 
The \textit{Vision Transformer (ViT)} treats image patches as tokens and feeds them into a transformer encoder to model global relationships between patches \cite{dosovitskiy2020image}.

For semantic segmentation tasks, transformer-based architectures typically follow a three-stage pipeline:

\begin{enumerate}
\item \textbf{Patch Embedding}

An input image $I \in \mathbb{R}^{H \times W \times C}$ is divided into non-overlapping patches of size $P \times P$. 
Each patch is flattened and projected into an embedding vector.

\item \textbf{Transformer Encoder}

The sequence of patch embeddings is processed by a transformer encoder consisting of multi-head self-attention layers and feed-forward networks. 
These layers model global dependencies across image regions.

\item \textbf{Segmentation Decoder}

A decoder module reconstructs pixel-level predictions from the encoded features. 
This stage usually includes upsampling layers or multi-scale feature fusion to produce the final segmentation map.
\end{enumerate}

Several transformer-based segmentation models have been proposed, including Vision Transformer (ViT), Swin Transformer, Segmenter, and SegFormer \cite{dosovitskiy2020image,liu2021swin,xie2021segformer,strudel2021segmenter}. 
These models demonstrate strong performance in dense prediction tasks due to their ability to capture global context.

\subsection{Datasets for Transformer-Based Segmentation}

Transformer-based segmentation models are commonly evaluated using large-scale annotated datasets. 
These datasets provide pixel-level annotations required for training and benchmarking segmentation models.

\textbf{Cityscapes}

The Cityscapes dataset is designed for urban scene understanding in autonomous driving applications \cite{cordts2016cityscapes}. 
It contains 5000 finely annotated images and 20000 coarsely labeled images with a resolution of $1024 \times 2048$. 
The dataset includes 19 semantic classes such as roads, vehicles, pedestrians, and traffic signs.

A key challenge in Cityscapes is the presence of small objects and complex boundaries, which require models to capture both local details and global context.

\textbf{ADE20K}

The ADE20K dataset is widely used for scene parsing and contains more than 25,000 images with annotations for 150 semantic classes \cite{zhou2017scene}. 
The dataset includes a diverse set of indoor and outdoor scenes, making it suitable for evaluating the generalization ability of segmentation models.

Transformer-based architectures perform well on ADE20K because the attention mechanism helps model relationships between multiple objects within complex scenes.

\textbf{PASCAL VOC 2012}

PASCAL VOC is a benchmark dataset commonly used in semantic segmentation research \cite{everingham2010pascal}. 
It includes approximately 20 object classes with pixel-level annotations.

Although smaller than other datasets, it is widely used for benchmarking segmentation algorithms.

\subsection{Evaluation Metrics}

To evaluate the performance of segmentation models, several pixel-level metrics are used.

\textbf{Pixel Accuracy (PA)}

Pixel accuracy measures the ratio of correctly predicted pixels to the total number of pixels:

\begin{equation}
PA = \frac{\sum_i n_{ii}}{\sum_i t_i}
\end{equation}

where $n_{ii}$ represents correctly classified pixels for class $i$, and $t_i$ represents the total number of pixels belonging to class $i$.

\textbf{Intersection over Union (IoU)}

Intersection over Union measures the overlap between the predicted segmentation and the ground truth:

\begin{equation}
IoU = \frac{TP}{TP + FP + FN}
\end{equation}

where $TP$ denotes true positives, $FP$ denotes false positives, and $FN$ denotes false negatives.

\textbf{Mean Intersection over Union (mIoU)}

Mean IoU is the average IoU across all classes:

\begin{equation}
mIoU = \frac{1}{K}\sum_{k=1}^{K} IoU_k
\end{equation}

where $K$ is the number of classes. 
This metric is considered the standard evaluation metric in most segmentation benchmarks.

\subsection{Research Gap}

Despite the success of transformer-based segmentation models, several limitations remain.

\begin{itemize}

\item \textbf{Computational Complexity}

Self-attention has quadratic complexity with respect to the number of tokens, which leads to high computational cost when processing high-resolution images.

\item \textbf{Large Data Requirement}

Transformer models typically require large-scale datasets to achieve optimal performance because they lack the inductive biases present in convolutional neural networks.

\item \textbf{Boundary Precision}

Although transformers capture global context effectively, they sometimes struggle to preserve fine spatial details, leading to inaccurate segmentation boundaries.

\item \textbf{Memory Consumption}

Large transformer models require significant GPU memory, which can limit their deployment in real-time applications or resource-constrained environments.

\end{itemize}

Future research directions include developing lightweight transformer architectures, improving data efficiency, and combining convolutional inductive biases with transformer attention mechanisms \cite{xie2021segformer,liu2021swin}.
